% =============================================================================
% 2.1 气象数据检索与分析的技术难点建模
% 草稿版本：v0.1 (2026-05-05)
% 字数：约 2200 字
% 风格规则：v0.2 风格（无 \textbf / \textit / 引例括号 / 双引号 / 破折号）
% 引用：\cite{ref2, ref5, ref10, ref20, ref25}
% 图表：
%   - 表 \reftab{tab:three-questions}     三个研究问题的对照建模
%   - 表 \reftab{tab:hallucination-types} 大语言模型在气象任务中的 5 类幻觉
% 依赖宏包：booktabs / tabularx（主稿已加载）
% =============================================================================

\subsection{气象数据检索与分析的技术难点建模}
\label{sec:problem-modeling}

把大语言模型应用到气象预报数据的检索与分析这一具体业务场景，本质上是
把通用的语言模型能力嫁接到一个对数值精度、来源权威性与时空一致性都有
严格要求的垂直领域。本节按用户感知到 Agent 内部的处理顺序，把这一嫁接
过程中暴露出的技术难点抽象为三个相互独立又相互衔接的研究问题，依次为
气象任务指令解析问题、气象数据异构特征问题、大语言模型数值与逻辑幻觉
问题，并在 \reftab{tab:three-questions} 中给出三者的对照建模。本章
后续内容将给出本文的统一应对方案，第 3 章至第 4 章则是这一方案在三类
问题上的具体落地。

\begin{table}[!ht]
    \centering
    \caption{气象 Agent 三类技术问题的对照建模}
    \label{tab:three-questions}
    \song\fontsize{9pt}{12pt}\selectfont
    \renewcommand{\arraystretch}{0.95}
    \begin{tabular}{lll}
        \toprule
        研究问题 & 主要表现 & 对应章节 \\
        \midrule
        指令解析问题 & 自然语言到结构化检索参数的映射缺口 & \ref{chap:retrieval} \\
        异构数据问题 & 多源、多粒度、多要素接入与归一化 & \ref{sec:api-tooling} 与 \ref{sec:semantic-bridge} \\
        数值与逻辑幻觉问题 & 数值理解、引用编造、统计口径漂移 & \ref{sec:semantic-bridge} 与 \ref{sec:code-execution} \\
        \bottomrule
    \end{tabular}
\end{table}

\subsubsection*{2.1.1 \quad 气象任务指令解析问题}

气象 Agent 接收的输入是用户自然语言查询，输出最终需要进入工具调用层
转化为带有结构化参数的 API 请求。两端之间存在显著的形式与语义鸿沟。

第一类形式鸿沟是参数的结构化要求。和风天气 \texttt{get\_\bk{}forcast\_\bk{}weather}
要求传入 \texttt{LocationID} 加 \texttt{hours} 取值集合中的一个枚举值
\texttt{24h}、\texttt{72h} 或 \texttt{168h}；Open-Meteo 历史数据接口
要求经度纬度加起止日期的 ISO 8601 字符串。用户却倾向于使用自由
表达式例如武汉这周天气如何或者去年这个月每小时降水多少。把武汉这周
解析为 \texttt{LocationID=\bk{}101200101} 加 \texttt{hours=\bk{}168h},把去年
这个月解析为 \texttt{[2024-04-01, 2024-04-30]}，需要 Agent 同时处理
命名实体识别、相对时间消歧、参数枚举映射三类子任务。代刊等 \cite{ref5}
在天气预报 LLM 应用综述中指出这类参数化转换是限制大语言模型在气象
业务上落地的首要瓶颈。

第二类语义鸿沟是用户表达的不完整性与角色多样性。用户提问明天天气怎么样
通常默认查询所在地，提问北京天气怎么样通常默认查询当天，提问明天早八
到晚六从武汉到成都有什么穿衣建议则隐含起点终点两个角色与对应时间段的
对齐关系。这类信息在前置 NLU 阶段如果不显式抽取，下游 ReAct 循环
会陷入猜参数的低效尝试。本文据此把指令解析问题拆解为意图识别加参数
补全两个相对独立的子任务，分别在 \ref{sec:intent-recognition} 节
与 \ref{sec:param-completion} 节展开。

第三类是对抗鲁棒性。用户输入并不总是合法气象查询，例如推荐部好看的
电影这类与气象无关的输入也会进入系统，需要 Agent 能识别并友好引导
而不是强行解析为某个气象意图。

\subsubsection*{2.1.2 \quad 气象数据的异构特征分析}

气象数据天然异构，且这一异构性贯穿数据接入、数据组织、数据解读三个
不同层面 \cite{ref25}。

第一是数据源异构。本文同时接入和风天气 API 与 Open-Meteo 两类数据源。
和风天气提供商业级实时观测、逐小时与逐日预报、天气预警、生活指数等
增值服务，但仅覆盖近 30 天与未来 30 天的窗口；Open-Meteo 由 European
Centre for Medium-Range Weather Forecasts 等机构维护，提供免费的
全球长期历史数据归档，时间深度可追溯到 1940 年代，但缺少天气预警与
生活指数等业务化能力。两类数据源在地理标识、参数命名、字段命名、
错误码体系上完全不一致，需要在工具层做统一封装。

第二是数据粒度异构。同一气象要素在不同接口下的时间粒度可能从分钟级
到月级跨度数个数量级，例如逐小时预报的 \texttt{precip} 字段以毫米
每小时为单位，逐日预报的 \texttt{precip} 字段以毫米每日为单位，历史
归档的 \texttt{precipitation\_\bk{}sum} 字段则是日累积量。直接把不同粒度
的数值传给大语言模型容易引发口径混乱例如把日累积 \SI{50}{mm} 当成
小时累积 \SI{50}{mm} 来判定。

第三是数据要素异构。气象要素中既有连续型如温度、湿度、降水、风速,
又有离散型如风力等级、能见度等级、生活指数等级，还有事件型如预警
信号、灾害等级。本文在 \ref{sec:semantic-bridge} 节通过双通路 RAG
架构把这一异构性吸收在确定性分类器内部，对外统一输出 \texttt{原始值
→ 分级名 → 影响描述 → 标准引用} 四元组。

需要特别指出的是，李特等 \cite{ref2} 在中国气象学会年会论文中提出的
气象 AI 应用思路明确把异构数据接入纳入 Agent 工程范畴，而非作为
传统数据中台问题处理。本文沿用这一思路，把异构性的解决方案分散到
工具学习与语义桥接两个 Agent 内部模块，避免在 Agent 之外维护额外
的数据中间层。

\subsubsection*{2.1.3 \quad 大语言模型的数值与逻辑幻觉}

大语言模型的训练目标本质上是以下一个 token 的预测来逼近文本概率
分布，这一目标天然导致模型在严格要求数值精度与权威引用的任务上
存在内在边界 \cite{ref10}。Hong 等 \cite{ref20} 在 Data Interpreter
中把这类问题概括为大语言模型在数据科学任务中的不可靠数值推理与
不可信生成两类失败。本文在气象领域的初步实证（\ref{sec:bridge-eval}
节与 \ref{sec:code-eval} 节）进一步把这一现象细化为五类典型幻觉,
归纳于 \reftab{tab:hallucination-types}。

\begin{table}[!ht]
    \centering
    \caption{大语言模型在气象任务中的 5 类典型幻觉}
    \label{tab:hallucination-types}
    \song\fontsize{9pt}{12pt}\selectfont
    \renewcommand{\arraystretch}{0.95}
    \begin{tabular}{lll}
        \toprule
        幻觉类型 & 主要表现 & 实证出处 \\
        \midrule
        引用编造 & 编造 GB/T 标准号或条款号 & \ref{sec:bridge-eval} 节 \\
        分级编造 & 自创不存在的等级编号或等级名 & \ref{sec:bridge-eval} 节 \\
        数值漂移 & 数数错、舍入不一致、字段命名漂移 & \ref{sec:code-eval} 节 \\
        范围错位 & 把日累积当小时值或反之 & \ref{sec:bridge-eval} 节 \\
        场景误判 & 不应输出建议时仍输出建议 & \ref{sec:bridge-eval} 节 \\
        \bottomrule
    \end{tabular}
\end{table}

引用编造是数值与逻辑幻觉问题中最典型也最危险的一类幻觉。在
\ref{sec:bridge-eval}
节的对照实验中，大语言模型自由发挥档对 71 条评测集输出的引用条款中
仅 5.6\% 的 \texttt{grade\_\bk{}id} 完全正确，且即使引用文本看起来合理,
\texttt{must\_\bk{}cite} 关键词的严格命中率也只有 25.4\%。换言之，模型
有约四分之三的概率会以看似权威的方式给出实际错误的引用，这种真实
却错的失败模式比纯粹捏造更难被用户识别。

分级编造是引用编造的姐妹失败模式。模型在被问到 \SI{35}{mm} 的降水
属于什么级别时会自创类似 \texttt{grading\_\bk{}precip\_\bk{}24h\_\bk{}strong} 这种
看起来合理但知识库中并不存在的 \texttt{grade\_\bk{}id}。这种自创编号在
形式上完全符合命名规则，足以骗过基于字符串匹配的下游引用检查。

数值漂移在 \ref{sec:code-eval} 节的代码执行评测中得到量化。在 18 条
覆盖求平均、求极值、求和、求差、计数、排序六类统计模式的评测集上,
让大语言模型直接做心算或自由生成结果，整体数值准确率为 66.7\%；
其中极值类任务准确率为 0\%，计数类任务准确率约 50\%。模型在需要维护
伴随状态的离散计数任务上失败率最高，这一现象与 Hong 等 \cite{ref20}
观察到的伴随状态不一致问题相吻合。

范围错位与场景误判是两类相对隐蔽但同样常见的幻觉。前者源于大语言
模型缺乏稳定的单位与时间粒度记忆，在不同上下文中对相同数值的解读
会发生切换；后者源于模型倾向于在结构化输出场景下也给出建议，例如
在户外作业场景下仍输出洗车提示，违反场景过滤约束。

针对数值与逻辑幻觉问题的五类幻觉，本文不寻求通过更大模型或更复杂
prompt 来缓解,
而是通过把数值理解、引用归属、统计计算三类高风险任务从大语言模型
内部转移到外部确定性组件。\ref{sec:semantic-bridge} 节用确定性分级
加 \texttt{grade\_\bk{}id} 硬链接消除引用与分级编造，
\ref{sec:code-execution} 节用代码生成加受限沙箱执行消除数值漂移。
本文据此把这一思路概括为大语言模型只做语言组织、数值与引用让外部
组件兜底，并以模块层与端到端两个维度的实证验证其有效性。

\subsubsection*{2.1.4 \quad 三个问题的层级关系}

三个研究问题在数据流向上构成自然的层级。指令解析问题解决用户输入
到结构化请求的转化，是 Agent 流水线的入口；异构数据问题解决结构化
请求到原始数据的取数与归一化，是数据通路的中段；数值与逻辑幻觉
问题解决原始数据到自然语言决策报告的解读与表达，是 Agent 流水线
的出口。三个问题的解决方案在工程上也存在依赖关系，异构数据接入
依赖指令解析输出的结构化意图作为参数来源，语义桥接依赖工具调用
结果作为输入数值。这种依赖关系决定了本文在第 3 章先解决前两个问题
对应的接入侧任务，第 4 章再解决数值与逻辑幻觉问题对应的解读侧任务。

作为本章后续内容的承接，下一节将讨论统一这三个子问题的方法论
框架，即基于 ReAct 范式的单智能体思维链模型。

